% !TEX root = ../main.tex

\begingroup
\makeatletter
\patchcmd{\@makechapterhead}{\vspace*{50\p@}}{\vspace*{0\p@}}{}{}
\makeatother
\chapter{Implementation}
\label{ch:implementation}
\endgroup

\section{Tools and language}
\label{sec:tools}

This design is implemented in VHDL 2008 and targets the Nexys A7-50T development board.
All signals are \lstinline|STD_LOGIC_VECTOR| and arithmetic operations (addition, subtraction,
multiplications, comparisons) are handled through the \lstinline|IEEE.NUMERIC_STD_UNSIGNED|
package (supported only in VHDL 2008). Vivado 2022.2 is used for simulation, synthesis and programming the board.

\section{Input classification and range reduction unit}
\label{sec:classification_rr}

The input classification and range reduction units are implemented as one entity because
both use the same unpacked fields (see \cref{sec:range-red-unit}).

\subsection{The entity}
\label{subsec:icrr-entity}
The unit takes the inputs \lstinline|x_in| and \lstinline|valid_in| and outputs the side information as an 8-bit \lstinline|STD_LOGIC_VECTOR|
rather than separate outputs for cleaner assignments in the pipeline (\cref{lst:classification_reduction_unit}), together with the
\lstinline|z_out| for the core, \lstinline|rsmall_out| for the small path and \lstinline|x_out|
for the bypass path. The \lstinline|valid_out| output bit marks which output cycles hold a real result.

\begin{lstlisting}[caption = {Input classification and range reduction entity declaration}, label = {lst:classification_reduction_unit}]
ENTITY classification_reduction_unit IS
    PORT (
        clk : IN STD_LOGIC; rst : IN STD_LOGIC; en : IN STD_LOGIC;
        x_in : IN STD_LOGIC_VECTOR(31 DOWNTO 0);
        valid_in : IN STD_LOGIC;

        side_info_out : OUT STD_LOGIC_VECTOR(7 DOWNTO 0);
        -- side_info_out bits significance:
        -- 7 = sign_r;
        -- 6 = sign_x;
        -- 5:4 = kmod4;
        -- 3:1 = input_cls
        -- 0 = small input flag
        x_out : OUT STD_LOGIC_VECTOR(31 DOWNTO 0);
        z_out : OUT STD_LOGIC_VECTOR(45 DOWNTO 0);      -- |r|, 2 int + 44 frac bits
        rsmall_out : OUT STD_LOGIC_VECTOR(47 DOWNTO 0); -- |r| bits 2^-13 .. 2^-60
        valid_out : OUT STD_LOGIC
    );
END classification_reduction_unit;
\end{lstlisting}

\subsection{Pipeline stages}
\label{subsec:icrr-stages}

\paragraph{Stage 1.}
The first stage classifies the input and gives a 3-bit encoding that covers the six
classes, as the case when $|x| \le \pi/4$ can be treated as a normal number, with $k = 0$
(\cref{lst:icrr-stage1}). Comparing the exponent with $115$ detects detects the tiny normal numbers
with $ |x| < 2^{-12}$ (the bias $127 - 12 = 115$, see \cref{sec:input-classification}, item 3).

The order of the \texttt{IF/ELSIF} blocks is important to check the inputs by priority.
Zero and subnormals have exponent $0$, which is also below 115.
If the \lstinline|x_exp < 115| test came first, they would be classified as tiny.
Exponents $255$ and $0$ are tested first, therefore the tiny test sees only normal numbers.

At the same time, the significand $M$ is prepared by adding the hidden $1$ bit back.

\begin{lstlisting}[caption = {First stage -- classification process logic}, label = {lst:icrr-stage1}]
-- unpack input
x_sign <= x_in(31);    x_exp <= x_in(30 DOWNTO 23);    x_frac <= x_in(22 DOWNTO 0);
-- process
IF x_exp = "11111111" THEN
    IF x_frac = 0 THEN stage1_class <= "100"; -- infinity
    ELSE               stage1_class <= "101"; -- NaN
    END IF;
ELSIF x_exp = "00000000" THEN
    IF x_frac = 0 THEN stage1_class <= "000"; -- zero
    ELSE               stage1_class <= "001"; -- subnormal
    END IF;
ELSIF x_exp < 115 THEN stage1_class <= "010"; -- tiny, |x| < 2^-12
ELSE                   stage1_class <= "111"; -- normal
END IF;
stage1_M <= '1' & x_frac;
stage1_sign_x <= x_sign;
\end{lstlisting}

\vspace{-6pt}
\paragraph{Stage 2.}
The second stage of the unit computes the window needed for the $x \cdot 2/\pi$ multiplication (\cref{lst:icrr-stage2}).
The windows are stored in an array of $256$ entries, with $82$ bits each, synthesized as a
Read Only Memory.

Each entry of the table has 82 bits: 2 integer bits, which give $k \bmod 4$, and 80
fractional bits. The number of fractional bits follows from the cancellation described in
\cref{sec:range-reduction}: when $x \cdot 2/\pi$ lies very close to an integer, the leading
bits of $r$ cancel out, and the remaining bits must still be correct. Over all binary32
inputs, the worst case gives $|r| \approx 2^{-29}$ quarter turns, i.e.\ 29 leading zeros.
Since the result needs about 27 correct bits after the first non-zero bit of $r$ (24 bits of
the significand and a few bits for rounding), $|r|$ must be accurate to about $2^{-56}$.
Truncating the window after the bit with weight $2^{-80}$ introduces an error below
$2^{-80}$, which is multiplied by the significand $M < 2^{24}$, giving a total error below
$2^{24} \cdot 2^{-80} = 2^{-56}$.

\begin{lstlisting}[caption = {Second stage -- window selection}, label = {lst:icrr-stage2}]
TYPE window_array IS ARRAY (0 TO 255) OF STD_LOGIC_VECTOR(81 DOWNTO 0);
CONSTANT WINDOW : window_array := (
    ...
    "0010100010111110011000...0011111110000100111", -- e = 150 (2/pi itself), 82 bits
    ...
);
...
stage2_W <= WINDOW(to_integer(stage1_exp));
stage2_sign_x <= stage1_sign_x;
\end{lstlisting}

\newpage

\paragraph{Stage 3.}
The third stage multiplies the significand $M$ by the selected window in the previous
stage, outputting $x \cdot 2/\pi$ (\cref{lst:icrr-stage3}).

This is done by multiplying the High and Low parts
in parallel, thus halving the critical path. Only the lower $82$ bits are kept,
because higher bits are multiples of $4$.

\begin{lstlisting}[caption = {Third stage -- multiplication}, label = {lst:icrr-stage3}]
-- stage 3a - parallel multiplication
stage3a_pp_hi <= stage2_M * stage2_W(81 DOWNTO 41);
stage3a_pp_lo <= stage2_M * stage2_W(40 DOWNTO 0);
stage3a_sign_x <= stage2_sign_x;

-- stage 3b - final product reconstruction
stage3b_prod <= (stage3a_pp_hi(40 DOWNTO 0) & (40 DOWNTO 0 => '0'))
              + ((16 DOWNTO 0 => '0') & stage3a_pp_lo);
stage3b_sign_x <= stage3a_sign_x;
\end{lstlisting}

\paragraph{Stage 4.}
The last stage of the input classification and range reduction unit takes the product
$P = x \cdot 2/\pi$ and rounds $k$ to the nearest integer depending on which multiple of
$\pi/2$ the product is closest to. If the first bit after the decimal point of $P$ (with
weight $0.5$) is \lstinline|'0'|, the nearest integer is the integer part itself
(bits $81$-$80$ of $P$). Otherwise, $k$ is rounded up .
On $2$ bits, $3+1$ wraps to $0$, so this is exactly mod $4$.

When rounding down, $|r|$ is the fraction and $s_r = 0$; when rounding up, $|r| = 1 - \text{fraction}$
(computed as \lstinline|0 - fraction| on 80 bits) and $s_r = 1$. If the first 12 fractional bits of
$|r|$ are zero, $|r| < 2^{-12}$ and the small-argument flag is set. Finally, the top 44 bits of
$|r|$ go to the core (\lstinline|z_out|), the bits $2^{-13}$ to $2^{-60}$ to the small-argument
path (\lstinline|rsmall_out|), and the remaining information to \lstinline|side_info_out|.

\begin{lstlisting}[caption = {Fourth stage -- rounding and outputs}, label = {lst:icrr-stage4}]
IF stage3b_prod(79) = '0' THEN       -- fraction < 0.5
    abs_r := stage3b_prod(79 DOWNTO 0);
    kmod4 := stage3b_prod(81 DOWNTO 80);
    v_sign_r := '0';
ELSE                                 -- fraction >= 0.5: round k up
    abs_r := 0 - stage3b_prod(79 DOWNTO 0);
    kmod4 := stage3b_prod(81 DOWNTO 80) + 1;
    v_sign_r := '1';
END IF;
...
z_out      <= "00" & abs_r(79 DOWNTO 36);
side_info_out <= v_sign_r & stage3b_sign_x & kmod4 & stage3b_class & small_flag;
rsmall_out <= abs_r(67 DOWNTO 20);
\end{lstlisting}